\section{Exact recombination of the four centered channels}
\label{sec:tpc20-recombination}

We begin with the divisor coefficient of the finite model.  Put
\begin{equation}
 \lambda'_R(u)
 =u\sum_{b\le R/u}
 \frac{\mu(ub)\mu(b)}{\varphi(ub)}
 \qquad(u\le R),
 \label{eq:tpc20-lambda-prime}
\end{equation}
and set \(\lambda'_R(u)=0\) for \(u>R\).  The divisor formula for
Ramanujan sums gives
\begin{equation}
 \Lambda_R(n)=\sum_{u\mid n}\lambda'_R(u).
 \label{eq:tpc20-model-divisor}
\end{equation}
The coefficient is squarefree supported.  For squarefree \(u\le R\),
\begin{equation}
 \lambda'_R(u)
 =\mu(u)\frac{u}{\varphi(u)}
 \sum_{\substack{b\le R/u\\(b,u)=1}}
 \frac{\mu(b)^2}{\varphi(b)}.
 \label{eq:tpc20-lambda-positive}
\end{equation}

Define
\begin{equation}
 a(u)=-\mu(u)\log u,\qquad
 b_R(u)=a(u)-\lambda'_R(u).
 \label{eq:tpc20-ab-coefficients}
\end{equation}
For squarefree \(u\), put
\begin{equation}
 w_R(u)
 =\log u+
 \one_{u\le R}\frac{u}{\varphi(u)}
 \sum_{\substack{b\le R/u\\(b,u)=1}}
 \frac{\mu(b)^2}{\varphi(b)}.
 \label{eq:tpc20-wR}
\end{equation}
Then \(w_R(u)>0\) and
\begin{equation}
 b_R(u)=-\mu(u)w_R(u),
 \label{eq:tpc20-bR-mobius}
\end{equation}
with both sides understood as zero off the squarefree integers.

Let \(U\ge R\) be an integer such that every target in the separated
packet under consideration satisfies
\begin{equation}
 1\le n_m(j)=mj+h\le U.
 \label{eq:tpc20-target-horizon}
\end{equation}
For \((u,mH)=1\), define
\begin{equation}
 C_{m,u}(j)=\one_{u\mid mj+h}-\frac1u,
 \label{eq:tpc20-common-centered-kernel}
\end{equation}
and define the finite prime-calibration defect
\begin{equation}
 \eps_U(m)
 =\sum_{\substack{u\le U\\(u,mH)=1}}\frac{a(u)}u
  -\rho_{P,H}(m).
 \label{eq:tpc20-calibration-defect}
\end{equation}

\begin{theorem}[One coefficient for all four channels]
\label{thm:tpc20-common-recombination}
Under the setup of \cref{sec:tpc20-interface}, in particular
\(dH\le R\), and under \cref{eq:tpc20-target-horizon}, one has
pointwise
\begin{equation}
 \boxed{
 Z_m(j)=\eps_U(m)+
 \sum_{\substack{u\le U\\(u,mH)=1}}
 b_R(u)C_{m,u}(j).}
 \label{eq:tpc20-common-recombination}
\end{equation}
Consequently, with \(\eps_i=\eps_U(m_i)\),
\begin{align}
 Z_{m_1}(j)Z_{m_2}(j)
 ={}&\eps_1\eps_2
 +\eps_1\sum_v b_R(v)C_{m_2,v}(j)
 +\eps_2\sum_u b_R(u)C_{m_1,u}(j)
 \notag\\
 &+\sum_{u,v}b_R(u)b_R(v)
 C_{m_1,u}(j)C_{m_2,v}(j),
 \label{eq:tpc20-four-channel-recombined}
\end{align}
where the last three sums carry the horizon and coprimality restrictions
from \eqref{eq:tpc20-common-recombination}.  Thus the product
\[
 P_{m_1}P_{m_2}-P_{m_1}M_{m_2}
 -M_{m_1}P_{m_2}+M_{m_1}M_{m_2}
\]
has a common coefficient and a common centered kernel before any
absolute value is taken.
\end{theorem}

\begin{proof}
The von Mangoldt divisor identity and
\eqref{eq:tpc20-target-horizon} give
\[
 \Lambda(n_m(j))
 =\sum_{\substack{u\le U\\u\mid n_m(j)}}a(u).
\]
Every divisor in this sum is coprime to \(mH\) by
\eqref{eq:tpc20-target-coprimality}.  Adding and subtracting the density
of each divisibility condition gives
\begin{equation}
 P_m(j)
 =\eps_U(m)+
 \sum_{\substack{u\le U\\(u,mH)=1}}a(u)C_{m,u}(j).
 \label{eq:tpc20-prime-common-kernel}
\end{equation}

Similarly, \eqref{eq:tpc20-model-divisor} and the exact conditional
mean \eqref{eq:tpc20-model-density} give
\begin{equation}
 M_m(j)
 =\sum_{\substack{u\le R\\(u,mH)=1}}
 \lambda'_R(u)C_{m,u}(j).
 \label{eq:tpc20-model-common-kernel}
\end{equation}
Indeed, averaging the divisor indicators over the primitive classes
shows that the sum of \(\lambda'_R(u)/u\) in
\eqref{eq:tpc20-model-common-kernel} is exactly
\(\rho_{M,H,R}(m)\).  Subtracting
\eqref{eq:tpc20-model-common-kernel} from
\eqref{eq:tpc20-prime-common-kernel} proves
\eqref{eq:tpc20-common-recombination}; multiplication proves
\eqref{eq:tpc20-four-channel-recombined}.
\end{proof}

\begin{remark}[What has and has not been matched]
\label{rem:tpc20-recombination-scope}
\Cref{thm:tpc20-common-recombination} is a finite, pointwise identity.
It constructs the common divisor coefficient which TPC-19 left as an
analytic requirement.  It is not a finite-window prime number theorem
in progressions and does not estimate the resulting two-row sum.  The
only scalar not absorbed into the common kernel is
\(\eps_U(m)\); it is controlled uniformly in
\cref{sec:tpc20-uniform-calibration}.
\end{remark}

For a periodic function of \(j\), let \(\mathbb E_H\) denote the
complete-period average conditional on \((j,H)=1\).

\begin{proposition}[Connected determinant zero mode]
\label{prop:tpc20-common-zero-mode}
Let \((u,m_1H)=(v,m_2H)=1\), put
\[
 g=(u,v),\qquad \Delta=m_1-m_2.
\]
Then
\begin{equation}
 \boxed{
 \mathbb E_H\!
 \left[C_{m_1,u}C_{m_2,v}\right]
 =\frac1{uv}\left(g\one_{g\mid\Delta}-1\right).}
 \label{eq:tpc20-common-zero-mode}
\end{equation}
In particular, every \(g=1\) zero mode in the recombined
\(PP-PM-MP+MM\) channel vanishes termwise.  Moreover,
\begin{equation}
 g\one_{g\mid\Delta}-1
 =\sum_{r=1}^{g-1}e_g(r\Delta).
 \label{eq:tpc20-determinant-spectrum}
\end{equation}
\end{proposition}

\begin{proof}
The two congruences for \(j\) are individually soluble and have
conditional densities \(1/u\) and \(1/v\).  Their joint density is
\(g/(uv)\) precisely when their reductions modulo \(g\) agree.
Primitivity makes this compatibility equivalent to \(g\mid\Delta\);
otherwise the joint density is zero.  Expanding the two centered
indicators proves \eqref{eq:tpc20-common-zero-mode}.  The complete
additive-character sum modulo \(g\), with its zero frequency removed,
gives \eqref{eq:tpc20-determinant-spectrum}.
\end{proof}

The common coefficient also identifies exactly where the M\"obius signs
remain.  On squarefree support write
\begin{equation}
 u=ga,\qquad v=gb,
 \qquad g,a,b\text{ pairwise coprime}.
 \label{eq:tpc20-squarefree-gab}
\end{equation}
Then
\begin{equation}
 b_R(ga)b_R(gb)
 =\mu(a)\mu(b)w_R(ga)w_R(gb).
 \label{eq:tpc20-common-sign-cancellation}
\end{equation}
The determinant content \(g\) is unsigned.  If the TPC-18 row signs
\(\mu(d_1)\mu(d_2)\) are also included, the admissibility conditions
imply \((a,d_1)=(b,d_2)=1\), and hence
\begin{equation}
 \mu(d_1)\mu(d_2)\mu(a)\mu(b)
 =\mu(d_1a)\mu(d_2b).
 \label{eq:tpc20-row-frequency-sign-fusion}
\end{equation}
No coprimality between \(d_1a\) and \(d_2b\) is asserted: \(d_1,d_2\)
may intersect, as may \(a,d_2\) and \(b,d_1\).
